\documentclass{article}
\usepackage{amsmath,graphicx,booktabs}
\begin{document}
\section{Introduction}\label{sec:intro}
This is plain text with \emph{emphasis}, some \textbf{bold words} and
inline math $a^2 + b^2 = c^2$ in the middle of a sentence.
See Section~\ref{sec:intro} and the work of~\cite{knuth}.

A second paragraph, longer, so that a selection can span words and
lines, and so that an edit lands in the middle of running text.
\subsection{Lists}
\begin{itemize}
  \item First item
  \item Second item with $x$
  \begin{enumerate}
    \item Nested one
    \item Nested two
  \end{enumerate}
  \item Third item
\end{itemize}
\begin{description}
  \item[Term] Its description.
\end{description}
\section{Mathematics}
\begin{equation}\label{eq:one}
  e^{i\pi} + 1 = 0
\end{equation}
\begin{align}
  a &= b + c \\
  d &= e \nonumber
\end{align}
\[ \int_0^1 x\,dx = \frac{1}{2} \]
\subsection{Floats}
\begin{figure}[htbp]
  \centering
  \includegraphics[width=0.3\linewidth]{example-image}
  \caption{A picture.}\label{fig:one}
\end{figure}
\begin{table}[htbp]
  \centering
  \begin{tabular}{ll}
    \toprule
    A & B \\
    \midrule
    1 & 2 \\
    \bottomrule
  \end{tabular}
  \caption{A table.}\label{tab:one}
\end{table}
\subsubsection{Last}
Text at the end, referring to Figure~\ref{fig:one} and Table~\ref{tab:one}.
\begin{thebibliography}{9}
\bibitem{knuth} D. E. Knuth, \emph{The \TeX book}, 1984.
\end{thebibliography}
\end{document}
